% Options for packages loaded elsewhere
\PassOptionsToPackage{unicode}{hyperref}
\PassOptionsToPackage{hyphens}{url}
\documentclass[
  10pt,
]{article}
\usepackage{xcolor}
\usepackage[margin=1in]{geometry}
\usepackage{amsmath,amssymb}
\setcounter{secnumdepth}{5}
\usepackage{iftex}
\ifPDFTeX
  \usepackage[T1]{fontenc}
  \usepackage[utf8]{inputenc}
  \usepackage{textcomp} % provide euro and other symbols
\else % if luatex or xetex
  \usepackage{unicode-math} % this also loads fontspec
  \defaultfontfeatures{Scale=MatchLowercase}
  \defaultfontfeatures[\rmfamily]{Ligatures=TeX,Scale=1}
\fi
\usepackage{lmodern}
\ifPDFTeX\else
  % xetex/luatex font selection
\fi
% Use upquote if available, for straight quotes in verbatim environments
\IfFileExists{upquote.sty}{\usepackage{upquote}}{}
\IfFileExists{microtype.sty}{% use microtype if available
  \usepackage[]{microtype}
  \UseMicrotypeSet[protrusion]{basicmath} % disable protrusion for tt fonts
}{}
\makeatletter
\@ifundefined{KOMAClassName}{% if non-KOMA class
  \IfFileExists{parskip.sty}{%
    \usepackage{parskip}
  }{% else
    \setlength{\parindent}{0pt}
    \setlength{\parskip}{6pt plus 2pt minus 1pt}}
}{% if KOMA class
  \KOMAoptions{parskip=half}}
\makeatother
\usepackage{longtable,booktabs,array}
\usepackage{calc} % for calculating minipage widths
% Correct order of tables after \paragraph or \subparagraph
\usepackage{etoolbox}
\makeatletter
\patchcmd\longtable{\par}{\if@noskipsec\mbox{}\fi\par}{}{}
\makeatother
% Allow footnotes in longtable head/foot
\IfFileExists{footnotehyper.sty}{\usepackage{footnotehyper}}{\usepackage{footnote}}
\makesavenoteenv{longtable}
\usepackage{graphicx}
\makeatletter
\newsavebox\pandoc@box
\newcommand*\pandocbounded[1]{% scales image to fit in text height/width
  \sbox\pandoc@box{#1}%
  \Gscale@div\@tempa{\textheight}{\dimexpr\ht\pandoc@box+\dp\pandoc@box\relax}%
  \Gscale@div\@tempb{\linewidth}{\wd\pandoc@box}%
  \ifdim\@tempb\p@<\@tempa\p@\let\@tempa\@tempb\fi% select the smaller of both
  \ifdim\@tempa\p@<\p@\scalebox{\@tempa}{\usebox\pandoc@box}%
  \else\usebox{\pandoc@box}%
  \fi%
}
% Set default figure placement to htbp
\def\fps@figure{htbp}
\makeatother
\setlength{\emergencystretch}{3em} % prevent overfull lines
\providecommand{\tightlist}{%
  \setlength{\itemsep}{0pt}\setlength{\parskip}{0pt}}
\usepackage{booktabs}
\usepackage{longtable}
\usepackage{array}
\usepackage{multirow}
\usepackage{wrapfig}
\usepackage{float}
\usepackage{colortbl}
\usepackage{pdflscape}
\usepackage{tabu}
\usepackage{threeparttable}
\usepackage{threeparttablex}
\usepackage[normalem]{ulem}
\usepackage{makecell}
\usepackage{xcolor}
\usepackage{bookmark}
\IfFileExists{xurl.sty}{\usepackage{xurl}}{} % add URL line breaks if available
\urlstyle{same}
\hypersetup{
  pdftitle={Response surface alignment measure},
  hidelinks,
  pdfcreator={LaTeX via pandoc}}

\title{Response surface alignment measure}
\author{}
\date{\vspace{-2.5em}}

\begin{document}
\maketitle

\section{Set-up}\label{set-up}

AOD at the pair of gridboxes at longitude -171.5625 and latitude 51.875
and longitude -169.6875 and latitude 51.875 (shown in green and blue
respectively in Figure 1) are considered, for July.

\begin{figure}
\centering
\pandocbounded{\includegraphics[keepaspectratio]{alignment-measure_files/figure-latex/map-of-selected-gridboxes-1.pdf}}
\caption{Selected gridboxes highlighted in green and blue.}
\end{figure}

\newpage

\section{Sensitivity analysis}\label{sensitivity-analysis}

Sensitivity analysis is not carried out.

\section{Emulator fit and validation}\label{emulator-fit-and-validation}

To fit the emulator, we assume a linear mean function of the form
\[\hspace{2.6cm}  m(\textbf{x}) = \beta_0 + \beta_1 x_1 + \beta_2 x_2 + \dots + \beta_px_p.\]

The Gauss (squared exponential) kernel is used for the covariance
function, maximum posterior mode estimation for the parameters estimates
(trend, range, SD and the nugget term), and the \texttt{RobustGaSP}
package to fit the emulators.

\begin{longtable}[]{@{}ccc@{}}
\caption{Estimated parameters}\tabularnewline
\toprule\noalign{}
Ranges, trends, variance and nugget & Output 1 & Output 2 \\
\midrule\noalign{}
\endfirsthead
\toprule\noalign{}
Ranges, trends, variance and nugget & Output 1 & Output 2 \\
\midrule\noalign{}
\endhead
\bottomrule\noalign{}
\endlastfoot
1 & 2e+02 & 2e+02 \\
2 & 4e+02 & 3e+02 \\
3 & 6e+08 & 4e+28 \\
4 & 3e+02 & 3e+02 \\
5 & 7e+08 & 2e+27 \\
6 & 4e+10 & 6e+27 \\
7 & 3e+01 & 3e+01 \\
8 & 6e+01 & 6e+01 \\
9 & 2e+02 & 2e+02 \\
10 & 7e+08 & 2e+21 \\
11 & 1e+02 & 1e+02 \\
12 & 2e+10 & 3e+35 \\
13 & 1e+08 & 3e+07 \\
14 & 9e+07 & 2e+25 \\
15 & 6e+01 & 6e+01 \\
16 & 4e+02 & 3e+02 \\
17 & 2e+02 & 2e+02 \\
18 & 3e+10 & 8e+25 \\
19 & 7e+01 & 7e+01 \\
20 & 3e+02 & 3e+02 \\
21 & 1e+02 & 1e+02 \\
22 & 2e+09 & 4e+24 \\
23 & 2e+08 & 3e+24 \\
24 & 1e+09 & 1e+29 \\
25 & 5e+02 & 5e+02 \\
26 & 2e+09 & 5e+24 \\
27 & 2e+02 & 2e+02 \\
28 & 2e+02 & 2e+02 \\
29 & 3e+09 & 1e+24 \\
30 & 2e+09 & 8e+25 \\
31 & 6e+08 & 4e+25 \\
32 & 8e+01 & 8e+01 \\
33 & 5e+08 & 2e+21 \\
34 & 6e+02 & 6e+02 \\
35 & 8e+01 & 7e+01 \\
36 & 7e+01 & 7e+01 \\
37 & 4e+02 & 4e+02 \\
38 & 9e+01 & 8e+01 \\
1 & 0e+00 & 0e+00 \\
2 & 0e+00 & 0e+00 \\
3 & 0e+00 & 0e+00 \\
4 & 0e+00 & 0e+00 \\
5 & 0e+00 & 0e+00 \\
6 & 0e+00 & 0e+00 \\
7 & 4e+00 & 4e+00 \\
8 & 0e+00 & 0e+00 \\
9 & 0e+00 & 0e+00 \\
10 & 0e+00 & 0e+00 \\
11 & 0e+00 & 0e+00 \\
12 & 0e+00 & 0e+00 \\
13 & 0e+00 & 0e+00 \\
14 & 0e+00 & 0e+00 \\
15 & 0e+00 & 0e+00 \\
16 & 0e+00 & 0e+00 \\
17 & 0e+00 & 0e+00 \\
18 & 0e+00 & 0e+00 \\
19 & 0e+00 & 0e+00 \\
20 & 0e+00 & 0e+00 \\
21 & 0e+00 & 0e+00 \\
22 & 0e+00 & 0e+00 \\
23 & 0e+00 & 0e+00 \\
24 & 0e+00 & 0e+00 \\
25 & 0e+00 & 0e+00 \\
26 & 0e+00 & 0e+00 \\
27 & 0e+00 & 0e+00 \\
28 & 0e+00 & 0e+00 \\
29 & 0e+00 & 0e+00 \\
30 & 0e+00 & 0e+00 \\
31 & 0e+00 & 0e+00 \\
32 & 0e+00 & 0e+00 \\
33 & 0e+00 & 0e+00 \\
34 & 0e+00 & 0e+00 \\
35 & 0e+00 & 0e+00 \\
36 & 0e+00 & 0e+00 \\
37 & 0e+00 & 0e+00 \\
1 & 3e+04 & 3e+04 \\
1 & 0e+00 & 0e+00 \\
\end{longtable}

For validation of the emulators fitted, as well as calculation of the
alignment measure, 10000 new LHC points are generated and predictions
made at these points. The mean width of the 95\% CI's for these
predictions are shown in Table 3, along with the mean width of the LOO
CV prediction 95\% CI's. Note that the range of output 1 is 0.6 and of
output 2 is 0.6.

\begin{longtable}[]{@{}
  >{\raggedright\arraybackslash}p{(\linewidth - 6\tabcolsep) * \real{0.0814}}
  >{\centering\arraybackslash}p{(\linewidth - 6\tabcolsep) * \real{0.3488}}
  >{\centering\arraybackslash}p{(\linewidth - 6\tabcolsep) * \real{0.3023}}
  >{\centering\arraybackslash}p{(\linewidth - 6\tabcolsep) * \real{0.2674}}@{}}
\caption{Validation checks}\tabularnewline
\toprule\noalign{}
\begin{minipage}[b]{\linewidth}\raggedright
Output
\end{minipage} & \begin{minipage}[b]{\linewidth}\centering
Prediction 95\% CI mean width
\end{minipage} & \begin{minipage}[b]{\linewidth}\centering
LOO CV 95\% CI mean width
\end{minipage} & \begin{minipage}[b]{\linewidth}\centering
\% validated in LOO CV
\end{minipage} \\
\midrule\noalign{}
\endfirsthead
\toprule\noalign{}
\begin{minipage}[b]{\linewidth}\raggedright
Output
\end{minipage} & \begin{minipage}[b]{\linewidth}\centering
Prediction 95\% CI mean width
\end{minipage} & \begin{minipage}[b]{\linewidth}\centering
LOO CV 95\% CI mean width
\end{minipage} & \begin{minipage}[b]{\linewidth}\centering
\% validated in LOO CV
\end{minipage} \\
\midrule\noalign{}
\endhead
\bottomrule\noalign{}
\endlastfoot
1 & 0.1 & 0.1 & 93 \\
2 & 0.1 & 0.11 & 94 \\
\end{longtable}

\begin{figure}
\includegraphics[width=0.5\linewidth]{alignment-measure_files/figure-latex/loo-cv-plots-0, figures-side-1} \includegraphics[width=0.5\linewidth]{alignment-measure_files/figure-latex/loo-cv-plots-0, figures-side-2} \caption{LOO CV validation plot for output 1 (left) and output 2 (right).}\label{fig:loo-cv-plots-0, figures-side}
\end{figure}

\newpage

\section{Check of by-hand
calculations}\label{check-of-by-hand-calculations}

\begin{figure}
\includegraphics[width=0.5\linewidth]{alignment-measure_files/figure-latex/by-hand-checks-1} \includegraphics[width=0.5\linewidth]{alignment-measure_files/figure-latex/by-hand-checks-2} \includegraphics[width=0.5\linewidth]{alignment-measure_files/figure-latex/by-hand-checks-3} \includegraphics[width=0.5\linewidth]{alignment-measure_files/figure-latex/by-hand-checks-4} \caption{The predictions found by-hand match well with those made by the RobustGaSP package (left), with the latter minus the former found to be essentially zero (right). Top: Output 1. Bottom: Output 2.}\label{fig:by-hand-checks}
\end{figure}

Ahead of performing by-hand calculations of the partial derivatives at
the testing points, by-hand calculations for the predictions made at
these points are plotted against the predictions returned by the
\texttt{RobustGaSP} package (left hand plots in Figure 3) and the
difference between the two plotted against the test input index (right
hand plots in Figure 3). Finally, for both we calculate
\[\dfrac{\sum\limits_{i=1}^{N}\Biggl|\text{(predictions from package)}_{i} - \text{(predictions by-hand)}_{i}\Biggr|}{N},\]
where \(N\) is the number of test points to get 0 for output 1 and 0 for
output 2.

\section{Calculation of the alignment
measure}\label{calculation-of-the-alignment-measure}

Having calculated the emulated partial derivatives for use in the
alignment measure, and normalising these, we calculate the AM, getting a
value of 1.

\end{document}
